\chapter{Generating schedule options}\label{ch:generate}

Once the inputs are valid, PlanWise can generate schedule options. Each algorithm yields one feasible
schedule from the same data; the system then keeps the schedules that no other schedule beats on every axis,
for you to compare (Chapter~\ref{ch:options}).

\shot{options_empty}{A period without options yet (here a what-if variant).}{options_empty}

\section{The Generate dialog}

Click \ui{Generate the first options} (or \ui{Generate options} / \ui{Generate more} in the period header).
You choose \emph{how long to wait}, not which solver:

\shotf[0.62\linewidth]{generate_dialog}{Choosing the search effort.}{generate_dialog}

\begin{center}
\begin{tabularx}{\linewidth}{>{\bfseries}p{2.2cm}p{3.6cm}X}
\toprule
\rowcolor{white}Preset & Runs & When to use \\
\midrule
Quick & EDD rule + Hybrid, 60~s & About a minute --- enough to see where things stand and how many orders are late. Good for a first look or a what-if variant. \\
Balanced & EDD rule + Hybrid, 180~s & Recommended for a two-week period; about three minutes. \\
Thorough & EDD rule + Hybrid, 600~s & Cuts operating cost further when there is time; about ten minutes. \\
\bottomrule
\end{tabularx}
\end{center}

The earliest-due-date rule (EDD) finishes in well under a second and is always included as a yardstick: an
optimised schedule must beat it to be worth the wait. The \term{Hybrid} method does the real work within the
budget, in three steps: dispatching rules give a fallback schedule; simulated annealing runs several seeds in
parallel to bring late orders down; then CP-SAT starts from the best schedule found to cut operating cost.
Each step replaces the previous result only when it is better, and it uses every processor core.

The footer shows how many runs will be queued and an upper bound of the waiting time. \ui{Generate
\textit{n} options} queues them.

\section{Customising algorithms, seeds and budget}

\ui{Customise algorithms, seeds and budget\dots} reveals every algorithm, grouped by family:

\shotf[0.6\linewidth]{generate_custom}{Custom selection of algorithms.}{generate_custom}

\begin{longtable}{>{\bfseries}p{3.5cm}p{2.4cm}p{8.6cm}}
\toprule
\rowcolor{white}Algorithm & Family & How it works \\
\midrule
\endhead
FIFO & Dispatching rule & First in, first out: at each step pick the operation that can start earliest; ties follow list order. \\
EDD & Dispatching rule & Earliest due date: as FIFO, ties favour the order due soonest. \\
SPT & Dispatching rule & Shortest processing time: as FIFO, ties favour the shortest run. \\
Simulated Annealing & Metaheuristic & Searches the priority order and machine choice, building schedules with the same decoder as the rules; accepts worse moves under a cooling temperature to escape local optima. \\
Genetic Algorithm & Metaheuristic & A population of solutions with tournament selection, crossover, mutation and elitism; schedules built with the same decoder. \\
CP-SAT & Constraint programming & One monolithic constraint model without hints: chooses machines, shifts, sequence and build-ahead options on its own. Needs a long budget on a two-week period. \\
CP-SAT + EDD hint & Constraint programming & CP-SAT warm-started from the EDD schedule; if it cannot improve, the EDD schedule is kept. \\
CP-LNS & Constraint programming & Large neighbourhood search: starts from the best rule schedule; each round reopens a time window, a few machines or one stage for CP-SAT to re-sequence. \\
Hybrid: search, then CP-SAT & Constraint programming & The three-step method described above; recommended for two-week periods. \\
Rolling-horizon CP-SAT & Constraint programming & Splits the horizon into consecutive windows, solves each with CP-SAT and freezes what is already scheduled. \\
\bottomrule
\end{longtable}

\begin{description}[style=nextline, leftmargin=1.2em, font=\bfseries\color{brand}]
  \item[Seeds for randomised algorithms] 1--3. The metaheuristics and CP-SAT searches are randomised; running
        two or three seeds gives a better chance and shows how stable the result is. Rules run once.
  \item[Budget per run (seconds)] The time limit of each non-rule run, up to the maximum published by the
        backend (900~s by default). A warning appears when CP-SAT gets less than two minutes: that is usually
        too short to improve a two-week schedule, so it will keep the EDD schedule.
\end{description}

\begin{note}
The dialog only offers algorithms the backend supports for the period's model. Hybrid and rolling-horizon
CP-SAT exist for the five-stage model only; for four-stage JSON scenarios (Advanced) the presets use rules,
simulated annealing, a genetic algorithm, CP-SAT with hint and CP-LNS instead.
\end{note}

\section{While the solver works}

Runs are queued and solved in the background on the server. You can keep working, open other pages or
close the browser. The period status becomes \ui{Solving}, the header shows ``Solving \textit{n} options
(Hybrid). The page updates when they finish.'', and the Options tab shows a spinner. When a run finishes, its
schedule is first checked by an \textbf{independent validator}; only valid schedules become options. The
page then updates on its own.

Each run ends with a solver status:

\begin{center}
\begin{tabularx}{\linewidth}{>{\bfseries}p{5.6cm}X}
\toprule
\rowcolor{white}Status & Meaning \\
\midrule
Feasible schedule & A valid schedule; not proven optimal (normal within a time limit). \\
Optimal & CP-SAT proved no better schedule exists for this model. \\
No improvement found, EDD schedule kept & A CP-SAT method ran out of time before improving the starting schedule. \\
No feasible schedule & The algorithm could not place everything; check the inputs (capacity, calendar). \\
\bottomrule
\end{tabularx}
\end{center}

\begin{tip}
If the analysis says ``The optimiser found nothing better than the dispatching rules'', it offers
\ui{Generate with Thorough}: the search and CP-SAT then get about ten minutes.
\end{tip}
